\section{Conclusions}
\label{sec:con}

\rev{We study how to allocate a fixed distillation budget across capability-targeted supervision to improve a specified target capability. ReAD combines a task-conditioned allocation prior, on-the-fly teacher supervision, and an uncertainty-aware contextual bandit. Under the same budget, it achieves higher on-target scores than target-only baselines, improves over static and greedy schedules with a matched generator, and shows smaller off-target declines than single-capability distillation. Its proxy reward may deviate from true task utility, and its efficiency figures are interpolation-based estimates from two budgets; validating both is left to future work. More broadly, our results indicate that how supervision is allocated across capabilities, and not only how much of it is used, shapes both what a distilled student learns and what it loses along the way.}
% demonstrating the value of treating capability distillation as a budgeted control problem that explicitly models capability interdependence.